% Source: https://artofproblemsolving.com/community/c3h2670574_some_actually_good_proving_problems
% Stable identity: aops-topic-2670574-q02
\emph{Lemma 1: $x\cdot0=0$.} By axioms 2 and 5, $x\cdot0=x(0+0)=x\cdot0+x\cdot0$. Adding $-(x\cdot0)$ to both sides and using associativity and axioms 2 and 3,
\[0=x\cdot0+\left(-(x\cdot0)\right)=\left(x\cdot0+x\cdot0\right)+\left(-(x\cdot0)\right)=x\cdot0+0=x\cdot0.\]

\emph{Lemma 2: if $u+v=0$ then $u=-v$.} Indeed $u=u+0=u+\left(v+(-v)\right)=(u+v)+(-v)=0+(-v)=-v$ (using commutativity for $0+(-v)=(-v)+0$).

\emph{Lemma 3: $(-1)(-1)=1$.} By axiom 5 and Lemma 1, $(-1)(-1)+(-1)\cdot1=(-1)\left((-1)+1\right)=(-1)\cdot0=0$, and $(-1)\cdot1=1\cdot(-1)=-1$ by axioms 1 and 4. So $(-1)(-1)+(-1)=0$ and Lemma 2 gives $(-1)(-1)=-(-1)$. Also $1+(-1)=0$, so Lemma 2 gives $1=-(-1)$. Hence $(-1)(-1)=1$.

\emph{Main proof.} By axiom 7, $1\ge0$ or $0\ge1$. If $0\ge1$, axiom 8 (adding $-1$) gives $0+(-1)\ge1+(-1)$, i.e. $-1\ge0$; then axiom 9 with $x=y=-1$ gives $(-1)(-1)\ge0$, i.e. $1\ge0$ by Lemma 3. So in every case $1\ge0$. Since $1\ne0$ (axiom 6), $1>0$ by definition. $\blacksquare$
